\documentclass[11pt,a4paper]{article}

\usepackage[utf8]{inputenc}
\usepackage[indonesian]{babel}
\usepackage{amsmath,amssymb,amsfonts,amsthm}
\usepackage{geometry}
\usepackage{tcolorbox}
\usepackage{booktabs}
\usepackage{enumitem}
\usepackage{hyperref}
\usepackage{tikz}
\usepackage{pgfplots}
\pgfplotsset{compat=1.18}

\geometry{
    top=2.5cm,
    bottom=2.5cm,
    left=2.5cm,
    right=2.5cm
}

\tcbuselibrary{skins,breakable}

% Colors
\definecolor{primaryblue}{RGB}{24, 76, 120}
\definecolor{accentblue}{RGB}{70, 130, 180}
\definecolor{solgreen}{RGB}{34, 112, 62}
\definecolor{solbg}{RGB}{246, 252, 248}
\definecolor{probblue}{RGB}{28, 70, 110}
\definecolor{probbg}{RGB}{245, 248, 253}

% Boxes for Problem and Solution
\newtcolorbox{problembox}[2][]{
    enhanced,
    breakable,
    colback=probbg,
    colframe=probblue,
    fonttitle=\bfseries\sffamily,
    title={#2 \hfill \normalfont\sffamily\footnotesize #1},
    boxrule=0.8pt,
    titlerule=0pt,
    coltitle=white,
    colbacktitle=probblue,
    arc=2.5mm,
    boxsep=2mm
}

\newtcolorbox{solutionbox}[1][]{
    enhanced,
    breakable,
    colback=solbg,
    colframe=solgreen,
    fonttitle=\bfseries\sffamily,
    title={Pembahasan #1},
    boxrule=0.7pt,
    titlerule=0pt,
    coltitle=white,
    colbacktitle=solgreen,
    arc=2mm,
    boxsep=2mm
}

\hypersetup{
    colorlinks=true,
    linkcolor=primaryblue,
    urlcolor=accentblue,
    citecolor=primaryblue
}

\title{\textbf{\LARGE Kumpulan Soal \& Pembahasan Lengkap Teori Suku Bunga}\\[0.3em]
\Large Modul 06: Anuitas dan Perpetuitas (Bagian 2)\\
\large (\textit{Unknown Time, Unknown Rate, Varying Rates, and Simple Interest})}
\author{\textbf{Standar:} SOA Exam FM / PAI A10 \& Perkuliahan Aktuaria ITB}
\date{\today}

\begin{document}

\maketitle
\noindent\rule{\textwidth}{0.4pt}
\vspace{1em}

\section*{Petunjuk Pengerjaan}
Kumpulan latihan soal ini mencakup topik lanjutan anuitas dan perpetuitas: penentuan waktu non-integer ($n+k$), pelunasan sisa saldo (*Balloon* vs *Drop Payments*), estimasi suku bunga tak diketahui (*Kellison Approximation* dan Interpolasi), evaluasi anuitas suku bunga bervariasi (*Portfolio Rate* vs *Yield Curve*), serta valuasi anuitas di bawah sistem bunga tunggal (*Simple Interest*).

\vspace{1em}

% ------------------- SOAL 1 -------------------
\begin{problembox}[Bobot: Dasar / Unknown Time \& Balloon vs Drop]{Soal 1}
Suatu pinjaman sebesar $\$10{,}000$ dikenakan suku bunga efektif tahunan $6\%$. Debitur melunasi pinjaman tersebut dengan angsuran tahunan sebesar $\$1{,}200$ di setiap akhir tahun selama mungkin, ditambah satu pembayaran sisa penyeimbang pada akhir masa pinjaman.
\begin{enumerate}[label=(\alph*)]
    \item Berapa kali pembayaran reguler penuh sebesar $\$1{,}200$ yang dapat dilakukan?
    \item Tentukan besaran pembayaran terakhir jika dibayarkan bersamaan dengan pembayaran reguler terakhir (\textit{Balloon Payment})!
    \item Tentukan besaran pembayaran terakhir jika dibayarkan satu tahun setelah pembayaran reguler terakhir (\textit{Drop Payment})!
\end{enumerate}
\end{problembox}

\begin{solutionbox}{Soal 1}
\textbf{Langkah 1: Menghitung Durasi Waktu $t$}
Persamaan nilai sekarang:
\begin{align*}
    10{,}000 &= 1{,}200 \cdot a_{\overline{t}|0.06} = 1{,}200 \left[ \frac{1 - (1.06)^{-t}}{0.06} \right] \\[0.4em]
    \frac{10{,}000 \times 0.06}{1{,}200} &= 0.50 = 1 - (1.06)^{-t} \implies (1.06)^{-t} = 0.50 \\[0.4em]
    (1.06)^t &= 2 \implies t = \frac{\ln(2)}{\ln(1.06)} = \frac{0.693147}{0.058269} \approx \mathbf{11.8957 \text{ tahun}}
\end{align*}
Jadi, terdapat **$11$ kali pembayaran reguler penuh** sebesar $\$1{,}200$ ($n = 11, k = 0.8957$).

\textbf{Langkah 2: Menghitung Sisa Saldo Pinjaman pada $t = 11$ ($B_{11}$)}
\begin{itemize}
    \item Akumulasi pinjaman pada $t = 11$: $10{,}000 \times (1.06)^{11} = 10{,}000 \times 1.898299 = \$18{,}982.99$
    \item Akumulasi 11 cicilan: $1{,}200 \times s_{\overline{11}|0.06} = 1{,}200 \times \left[\frac{1.898299 - 1}{0.06}\right] = 1{,}200 \times 14.97164 = \$17{,}965.97$
    \item Sisa saldo pinjaman persis setelah pembayaran ke-11:
    \begin{equation*}
        B_{11} = \$18{,}982.99 - \$17{,}965.97 = \mathbf{\$1{,}017.02}
    \end{equation*}
\end{itemize}

\textbf{Langkah 3: Menentukan Jenis Pembayaran Penyeimbang}
\begin{enumerate}[label=(\alph*)]
    \item Pembayaran penuh dapat dilakukan sebanyak **$11$ kali**.
    \item \textbf{\textit{Balloon Payment} pada $t = 11$:}
    \begin{equation*}
        \text{Total Pembayaran ke-11} = \$1{,}200 + B_{11} = \$1{,}200 + \$1{,}017.02 = \mathbf{\$2{,}217.02}
    \end{equation*}
    \item \textbf{\textit{Drop Payment} pada $t = 12$ ($X$):}
    \begin{equation*}
        X = B_{11} \times (1 + 0.06) = \$1{,}017.02 \times 1.06 = \mathbf{\$1{,}078.04}
    \end{equation*}
\end{enumerate}
\end{solutionbox}

\vspace{1.5em}

% ------------------- SOAL 2 -------------------
\begin{problembox}[Bobot: Standar SOA FM / Pembayaran Fraksional $t = n+k$]{Soal 2}
Mengacu pada transaksi pinjaman pada Soal 1 ($PV = \$10{,}000, R = \$1{,}200, i = 6\%, t \approx 11.8957$):
\begin{enumerate}[label=(\alph*)]
    \item Tentukan besaran pembayaran penyeimbang $X$ jika dibayarkan tepat pada waktu fraksional $t = 11.8957$!
    \item Tunjukkan bahwa nilai pembayaran sisa tersebut mendekati $k \cdot R = 0.8957 \times \$1{,}200$!
\end{enumerate}
\end{problembox}

\begin{solutionbox}{Soal 2}
\textbf{Langkah 1: Menghitung Pembayaran pada Waktu Fraksional $t = 11.8957$}
Sisa saldo $B_{11} = \$1{,}017.02$ diakumulasikan selama $k = 0.8957$ tahun pada suku bunga $6\%$:
\begin{align*}
    X_{11.8957} &= B_{11} \cdot (1.06)^{0.8957} = \$1{,}017.02 \times (1.053538) = \mathbf{\$1{,}071.47}
\end{align*}

\textbf{Langkah 2: Perbandingan dengan Pendekatan Linier ($k \cdot R$)}
Berdasarkan aproksimasi $a_{\overline{n+k}|} \approx a_{\overline{n}|} + k v^{n+k}$, sisa pembayaran proporsional terhadap cicilan $R$:
\begin{equation*}
    k \cdot R = 0.8957 \times \$1{,}200 = \mathbf{\$1{,}074.84}
\end{equation*}
Selisih antara nilai eksak $(\$1{,}071.47)$ dan estimasi linier $(\$1{,}074.84)$ hanya sebesar $\$3.37$ ($< 0.3\%$).
\end{solutionbox}

\vspace{1.5em}

% ------------------- SOAL 3 -------------------
\begin{problembox}[Bobot: Menengah / Unknown Rate Aproksimasi Aljabar]{Soal 3}
Suatu anuitas 12 tahun membayarkan $\$500$ di setiap akhir tahun dan memiliki nilai sekarang sebesar $\$4{,}200$.
\begin{enumerate}[label=(\alph*)]
    \item Tentukan estimasi suku bunga efektif tahunan $i$ menggunakan rumus aproksimasi Kellison orde-2!
    \item Perbaiki hasil estimasi tersebut menggunakan satu kali langkah iterasi Newton-Raphson!
\end{enumerate}
\end{problembox}

\begin{solutionbox}{Soal 3}
\textbf{Langkah 1: Identifikasi Nilai $K$ dan $n$}
\begin{itemize}
    \item $PV = \$4{,}200, R = \$500, n = 12$.
    \item $K = a_{\overline{12}|i} = \frac{4{,}200}{500} = 8.40$.
\end{itemize}

\textbf{Langkah 2: Estimasi Awal dengan Rumus Aproksimasi Kellison}
\begin{align*}
    i_0 &\approx \frac{2(n - K)}{K(n + 1)} = \frac{2(12 - 8.40)}{8.40(12 + 1)} = \frac{2(3.60)}{8.40 \times 13} = \frac{7.20}{109.20} \approx \mathbf{0.06593 \ (6.593\%)}
\end{align*}

\textbf{Langkah 3: Iterasi Newton-Raphson}
Definisikan $f(i) = \frac{1 - (1+i)^{-12}}{i} - 8.40$.
Untuk $i_0 = 0.06593$:
\begin{itemize}
    \item $v = (1.06593)^{-1} \approx 0.938148$
    \item $(1.06593)^{-12} \approx 0.463660 \implies a_{\overline{12}|0.06593} = \frac{1 - 0.463660}{0.06593} \approx 8.1350$
    \item $f(i_0) = 8.1350 - 8.40 = -0.2650$
    \item $f'(i_0) = -\frac{a_{\overline{12}|i_0} - 12 v^{13}}{i_0} = -\frac{8.1350 - 12(0.434983)}{0.06593} = -\frac{8.1350 - 5.2198}{0.06593} = -\frac{2.9152}{0.06593} \approx -44.217$
\end{itemize}
Pembaruan nilai $i$:
\begin{equation*}
    i_1 = i_0 - \frac{f(i_0)}{f'(i_0)} = 0.06593 - \frac{-0.2650}{-44.217} = 0.06593 - 0.00599 = \mathbf{0.05994 \ (5.994\%)}
\end{equation*}
(Nilai eksak adalah $i \approx 6.00\%$).
\end{solutionbox}

\vspace{1.5em}

% ------------------- SOAL 4 -------------------
\begin{problembox}[Bobot: Standar SOA FM / Portfolio Rate Method]{Soal 4}
Sebuah rekening tabungan berjangka 6 tahun menerima setoran sebesar $\$2{,}000$ di setiap akhir tahun ($t = 1, 2, \dots, 6$). Suku bunga efektif tahunan yang berlaku pada rekening tersebut berfluktuasi sebagai berikut:
\begin{itemize}
    \item Tahun ke-1 dan ke-2: $i_1 = i_2 = 4\%$
    \item Tahun ke-3 dan ke-4: $i_3 = i_4 = 5\%$
    \item Tahun ke-5 dan ke-6: $i_5 = i_6 = 6\%$
\end{itemize}
Hitunglah:
\begin{enumerate}[label=(\alph*)]
    \item Nilai terakumulasi saldo rekening pada akhir tahun ke-6 ($AV_6$) menggunakan \textit{Portfolio Rate Method}!
    \item Nilai sekarang dari seluruh setoran tersebut pada waktu $t = 0$ ($PV_0$)!
\end{enumerate}
\end{problembox}

\begin{solutionbox}{Soal 4}
\textbf{Langkah 1: Menghitung Nilai Akumulasi pada $t = 6$ ($AV_6$)}
Gunakan pergerakan saldo per periode waktu:
\begin{itemize}
    \item Akhir tahun ke-2 (setelah setoran ke-2):
    \begin{equation*}
        AV_2 = 2{,}000 \cdot s_{\overline{2}|0.04} = 2{,}000 \times 2.04 = \$4{,}080.00
    \end{equation*}
    \item Akhir tahun ke-4 (setelah setoran ke-4 pada bunga $5\%$):
    \begin{align*}
        AV_4 &= AV_2 \cdot (1.05)^2 + 2{,}000 \cdot s_{\overline{2}|0.05} \\
             &= 4{,}080 \times 1.1025 + 2{,}000 \times 2.05 = 4{,}498.20 + 4{,}100.00 = \$8{,}598.20
    \end{align*}
    \item Akhir tahun ke-6 (setelah setoran ke-6 pada bunga $6\%$):
    \begin{align*}
        AV_6 &= AV_4 \cdot (1.06)^2 + 2{,}000 \cdot s_{\overline{2}|0.06} \\
             &= 8{,}598.20 \times 1.1236 + 2{,}000 \times 2.06 = 9{,}660.94 + 4{,}120.00 = \mathbf{\$13{,}780.94}
    \end{align*}
\end{itemize}

\textbf{Langkah 2: Menghitung Nilai Sekarang pada $t = 0$ ($PV_0$)}
Diskontokan kembali seluruh akumulasi $AV_6$ dari $t = 6$ ke $t = 0$:
\begin{align*}
    PV_0 &= AV_6 \cdot (1.06)^{-2} \cdot (1.05)^{-2} \cdot (1.04)^{-2} \\
         &= 13{,}780.94 \times (0.889996) \times (0.907029) \times (0.924556) = \mathbf{\$10{,}287.41}
\end{align*}
\end{solutionbox}

\vspace{1.5em}

% ------------------- SOAL 5 -------------------
\begin{problembox}[Bobot: Standar SOA FM / Yield Curve Method vs Portfolio Method]{Soal 5}
Diberikan struktur suku bunga \textit{Spot Rate} ($s_t$) untuk periode $t = 1, 2, 3$ tahun:
\begin{equation*}
    s_1 = 4.0\%, \quad s_2 = 4.8\%, \quad s_3 = 5.5\%
\end{equation*}
\begin{enumerate}[label=(\alph*)]
    \item Tentukan nilai sekarang dari anuitas-immediate 3 tahun yang membayarkan $\$1{,}000$ per tahun menggunakan \textit{Yield Curve Method}!
    \item Tentukan suku bunga efektif tahunan konstan $i_{\text{level}}$ yang menghasilkan nilai sekarang yang ekuivalen!
\end{enumerate}
\end{problembox}

\begin{solutionbox}{Soal 5}
\textbf{Langkah 1: Menghitung Nilai Sekarang dengan Yield Curve (Spot Rates)}
Setiap arus kas didiskontokan dengan spot rate masing-masing tenor:
\begin{align*}
    PV &= 1{,}000 \cdot (1 + s_1)^{-1} + 1{,}000 \cdot (1 + s_2)^{-2} + 1{,}000 \cdot (1 + s_3)^{-3} \\[0.4em]
       &= 1{,}000 \left[ (1.040)^{-1} + (1.048)^{-2} + (1.055)^{-3} \right] \\[0.4em]
       &= 1{,}000 \left[ 0.961538 + 0.910491 + 0.851614 \right] \\[0.4em]
       &= 1{,}000 \times 2.723643 = \mathbf{\$2{,}723.64}
\end{align*}

\textbf{Langkah 2: Menentukan Suku Bunga Ekuivalen ($i_{\text{level}}$)}
Persamaan: $a_{\overline{3}|i_{\text{level}}} = 2.723643$.
\begin{itemize}
    \item Uji $i = 5.0\%$: $a_{\overline{3}|0.05} = \frac{1 - (1.05)^{-3}}{0.05} \approx 2.723248 \approx 2.723643$.
\end{itemize}
Suku bunga ekuivalen rata-rata adalah $\mathbf{i_{\text{level}} \approx 4.993\% \ (\approx 5.0\%)}$.
\end{solutionbox}

\vspace{1.5em}

% ------------------- SOAL 6 -------------------
\begin{problembox}[Bobot: Menengah / Anuitas Bunga Tunggal]{Soal 6}
Sebuah perjanjian kerja sama bisnis menetapkan pembayaran sewa sebesar $\$2{,}500$ di setiap akhir tahun selama 4 tahun ($t = 1, 2, 3, 4$). Perjanjian tersebut memberlakukan suku bunga tunggal (\textit{simple interest}) sebesar $8\%$ per tahun.
\begin{enumerate}[label=(\alph*)]
    \item Hitunglah nilai akumulasi ($AV$) dari total pembayaran sewa pada akhir tahun ke-4!
    \item Hitunglah nilai sekarang ($PV$) dari total pembayaran sewa pada awal tahun ke-1 ($t = 0$)!
\end{enumerate}
\end{problembox}

\begin{solutionbox}{Soal 6}
\textbf{Langkah 1: Menghitung Nilai Akumulasi ($AV$)}
Gunakan rumus $s_{\overline{n}|} = n + i \frac{n(n-1)}{2}$ untuk $n = 4, i = 0.08$:
\begin{align*}
    s_{\overline{4}|} &= 4 + 0.08 \times \frac{4 \times 3}{2} = 4 + 0.08 \times 6 = 4 + 0.48 = 4.48 \\[0.4em]
    AV &= 2{,}500 \times 4.48 = \mathbf{\$11{,}200.00}
\end{align*}

\textbf{Langkah 2: Menghitung Nilai Sekarang ($PV$)}
\begin{align*}
    PV &= 2{,}500 \sum_{t=1}^4 \frac{1}{1 + 0.08t} = 2{,}500 \left[ \frac{1}{1.08} + \frac{1}{1.16} + \frac{1}{1.24} + \frac{1}{1.32} \right] \\[0.4em]
       &= 2{,}500 \left[ 0.925926 + 0.862069 + 0.806452 + 0.757576 \right] \\[0.4em]
       &= 2{,}500 \times 3.352023 = \mathbf{\$8{,}380.06}
\end{align*}
\end{solutionbox}

\vspace{1.5em}

% ------------------- SOAL 7 -------------------
\begin{problembox}[Bobot: Menantang / SOA Exam FM Unknown Time]{Soal 7}
Sebuah rekening pensiun memiliki saldo awal $\$50{,}000$ dan menghasilkan suku bunga efektif tahunan $7\%$. Pemilik rekening berencana menarik dana sebesar $\$4{,}500$ di setiap akhir tahun.
\begin{enumerate}[label=(\alph*)]
    \item Berapa tahun pemilik rekening dapat melakukan penarikan penuh sebesar $\$4{,}500$?
    \item Jika sisa saldo terakhir ditarik sebagai \textit{Drop Payment} pada tahun ke-$(n+1)$, berapakah besar penarikan terakhir tersebut?
    \item Berapakah besar penarikan tahunan maksimum yang dapat ditarik selamanya tanpa pernah menghabiskan pokok tabungan (\textit{Perpetuity})?
\end{enumerate}
\end{problembox}

\begin{solutionbox}{Soal 7}
\textbf{Langkah 1: Menghitung Durasi Waktu $t$}
\begin{align*}
    50{,}000 &= 4{,}500 \cdot a_{\overline{t}|0.07} = 4{,}500 \left[\frac{1 - (1.07)^{-t}}{0.07}\right] \\[0.4em]
    \frac{50{,}000 \times 0.07}{4{,}500} &= \frac{3{,}500}{4{,}500} = 0.777778 = 1 - (1.07)^{-t} \implies (1.07)^{-t} = 0.222222 \\[0.4em]
    -t \ln(1.07) &= \ln(0.222222) \implies t = -\frac{-1.504077}{0.067659} \approx \mathbf{22.23 \text{ tahun}}
\end{align*}
Penarikan penuh $\$4{,}500$ dapat dilakukan sebanyak **$22$ kali** ($n = 22$).

\textbf{Langkah 2: Menghitung Saldo $B_{22}$ \& Drop Payment pada $t = 23$}
\begin{itemize}
    \item Saldo pada $t = 22$:
    \begin{align*}
        B_{22} &= 50{,}000(1.07)^{22} - 4{,}500 \cdot s_{\overline{22}|0.07} \\[0.3em]
        (1.07)^{22} &\approx 4.430402 \implies s_{\overline{22}|0.07} = \frac{3.430402}{0.07} \approx 49.00574 \\[0.3em]
        B_{22} &= 50{,}000(4.430402) - 4{,}500(49.00574) \\
               &= 221{,}520.10 - 220{,}525.83 = \mathbf{\$994.27}
    \end{align*}
    \item \textit{Drop Payment} pada $t = 23$:
    \begin{equation*}
        X = B_{22} \times 1.07 = \$994.27 \times 1.07 = \mathbf{\$1{,}063.87}
    \end{equation*}
\end{itemize}

\textbf{Langkah 3: Penarikan Maksimum Abadi (Perpetuitas)}
Agar pokok $\$50{,}000$ tidak pernah berkurang, penarikan tahunan sama dengan bunga yang dihasilkan setiap tahun:
\begin{equation*}
    R_{\text{perpetual}} = \$50{,}000 \times 7\% = \mathbf{\$3{,}500 \text{ per tahun}}
\end{equation*}
\end{solutionbox}

\vspace{1.5em}

% ------------------- SOAL 8 -------------------
\begin{problembox}[Bobot: Menantang / Suku Bunga Berpola Barisan]{Soal 8}
Suatu dana investasi membayarkan suku bunga pada tahun ke-$k$ sebesar $i_k = \frac{1}{k+1}$ untuk $k = 1, 2, 3, \dots, n$.
Tentukan rumus sederhana tertutup untuk nilai sekarang dari anuitas-immediate $n$-periode dengan pembayaran $\$1$ per tahun ($a_{\overline{n}|}$) menggunakan \textit{Portfolio Rate Method}!
\end{problembox}

\begin{solutionbox}{Soal 8}
\textbf{Langkah 1: Menghitung Faktor Diskon Akumulatif untuk Periode ke-$t$}
Suku bunga pada tahun ke-$s$ adalah $i_s = \frac{1}{s+1}$. Maka faktor pertumbuhan satu periode:
\begin{equation*}
    1 + i_s = 1 + \frac{1}{s+1} = \frac{s+2}{s+1}
\end{equation*}
Faktor diskonto dari waktu $t$ ke $t=0$:
\begin{equation*}
    \prod_{s=1}^t (1 + i_s)^{-1} = \prod_{s=1}^t \left( \frac{s+1}{s+2} \right) = \left(\frac{2}{3}\right) \left(\frac{3}{4}\right) \left(\frac{4}{5}\right) \dots \left(\frac{t+1}{t+2}\right) = \frac{2}{t+2}
\end{equation*}

\textbf{Langkah 2: Menjumlahkan Nilai Sekarang Seluruh Pembayaran}
\begin{align*}
    a_{\overline{n}|} &= \sum_{t=1}^n \prod_{s=1}^t (1+i_s)^{-1} = \sum_{t=1}^n \frac{2}{t+2} = 2 \sum_{t=1}^n \frac{1}{t+2} \\[0.4em]
                      &= 2 \left[ \frac{1}{3} + \frac{1}{4} + \frac{1}{5} + \dots + \frac{1}{n+2} \right]
\end{align*}
Untuk $n = 3$, $a_{\overline{3}|} = 2\left(\frac{1}{3} + \frac{1}{4} + \frac{1}{5}\right) = 2\left(\frac{20 + 15 + 12}{60}\right) = 2\left(\frac{47}{60}\right) = \frac{47}{30} \approx \mathbf{1.5667}$.
\end{solutionbox}

\vspace{1.5em}

% ------------------- SOAL 9 -------------------
\begin{problembox}[Bobot: Menengah / Komparasi Kinerja Sinking Fund]{Soal 9}
Perusahaan XYZ meminjam $\$100{,}000$ selama 5 tahun dengan bunga pinjaman $8\%$ per tahun yang dibayarkan di setiap akhir tahun. Untuk melunasi pokok pinjaman $\$100{,}000$ pada akhir tahun ke-5, perusahaan menyetorkan dana cadangan (*sinking fund*) di setiap akhir tahun ke rekening terpisah yang menghasilkan bunga $6\%$ per tahun.
\begin{enumerate}[label=(\alph*)]
    \item Tentukan total pengeluaran tahunan perusahaan XYZ!
    \item Bandingkan dengan cicilan tahunan jika pinjaman dilunasi melalui amortisasi langsung pada suku bunga $8\%$!
\end{enumerate}
\end{problembox}

\begin{solutionbox}{Soal 9}
\textbf{Langkah 1: Menghitung Pengeluaran Tahunan Metode Sinking Fund}
\begin{itemize}
    \item Bunga pinjaman tahunan: $I = \$100{,}000 \times 8\% = \$8{,}000$.
    \item Setoran \textit{sinking fund} tahunan ($D$):
    \begin{equation*}
        D \cdot s_{\overline{5}|0.06} = 100{,}000 \implies D = \frac{100{,}000}{\frac{(1.06)^5 - 1}{0.06}} = \frac{100{,}000}{5.637093} = \$17{,}739.65
    \end{equation*}
    \item Total pengeluaran tahunan $= \$8{,}000 + \$17{,}739.65 = \mathbf{\$25{,}739.65}$
\end{itemize}

\textbf{Langkah 2: Menghitung Cicilan Metode Amortisasi Langsung pada $8\%$}
\begin{equation*}
    R = \frac{100{,}000}{a_{\overline{5}|0.08}} = \frac{100{,}000}{\frac{1 - (1.08)^{-5}}{0.08}} = \frac{100{,}000}{3.992710} = \mathbf{\$25{,}045.65}
\end{equation*}
Metode amortisasi langsung lebih hemat $\$694.00$ per tahun karena dana pelunasan tidak terkena penalti selisih suku bunga ($8\% - 6\% = 2\%$).
\end{solutionbox}

\vspace{1.5em}

% ------------------- SOAL 10 -------------------
\begin{problembox}[Bobot: Standar Aplikasi Bisnis / Keputusan Investasi KPR]{Soal 10}
Seseorang membeli rumah seharga $\$200{,}000$ dengan membayar uang muka (*down payment*) $20\%$ dan sisanya dibiayai KPR dengan suku bunga $i^{(12)} = 7.2\%$ selama 20 tahun.
\begin{enumerate}[label=(\alph*)]
    \item Hitung cicilan bulanan KPR tersebut!
    \item Jika setelah 5 tahun debitur ingin melunasi seluruh sisa pinjaman sekaligus, berapa total dana pelunasan pokok yang harus dibayarkan?
\end{enumerate}
\end{problembox}

\begin{solutionbox}{Soal 10}
\textbf{Langkah 1: Menghitung Pokok Pinjaman \& Suku Bunga Bulanan}
\begin{itemize}
    \item Uang muka $= 20\% \times \$200{,}000 = \$40{,}000$.
    \item Pokok pinjaman KPR: $L = \$200{,}000 - \$40{,}000 = \$160{,}000$.
    \item Suku bunga bulanan: $j = \frac{7.2\%}{12} = 0.6\% = 0.006$.
    \item Total periode: $n = 20 \times 12 = 240 \text{ bulan}$.
\end{itemize}

\textbf{Langkah 2: Menghitung Cicilan Bulanan ($R$)}
\begin{align*}
    R &= \frac{L}{a_{\overline{240}|0.006}} = \frac{160{,}000}{\frac{1 - (1.006)^{-240}}{0.006}} = \frac{160{,}000}{\frac{1 - 0.237885}{0.006}} = \frac{160{,}000}{127.0191} = \mathbf{\$1{,}259.65 \text{ per bulan}}
\end{align*}

\textbf{Langkah 3: Menghitung Sisa Saldo Pelunasan Setelah 5 Tahun (Bulan ke-60)}
Sisa masa pinjaman yang belum berjalan adalah $240 - 60 = 180 \text{ bulan}$.
Gunakan metode prospektif:
\begin{align*}
    B_{60} &= R \cdot a_{\overline{180}|0.006} = 1{,}259.65 \left[ \frac{1 - (1.006)^{-180}}{0.006} \right] \\[0.4em]
           &= 1{,}259.65 \left[ \frac{1 - 0.340656}{0.006} \right] = 1{,}259.65 \times 109.8907 = \mathbf{\$138{,}423.86}
\end{align*}
Total dana yang harus disiapkan untuk melunasi pokok pada akhir tahun ke-5 adalah **$\$138{,}423.86$**.
\end{solutionbox}

\end{document}
